% Figure 7 draft D (merge of B and C): sites on the torsion circle with the mirror line and the t arc; the grid of
% disc pictures, each mode with its t image and b image; beneath, the fixed point of A1 and the E plane with the
% three images of v2 and the mirror axis. Labels only, one line on how to read the discs.
\documentclass[10pt,border=1mm]{standalone}
\usepackage[T1]{fontenc}
\usepackage{figurestyle}
\input{primitives.tex}
\input{numbers.tex}
% shared drawing for the figure 7 drafts: a mode picture is three discs on the three version sites;
% disc area = |amplitude|, ink positive, white negative, a dash for zero. Sites at 90, 210, 330 degrees.
% \modepic{x}{y}{R}{a}{b}{c}{style}  style: circle (torsion circle) or tri (triangle)
\newcommand\modepic[7]{\begin{scope}[shift={(#1,#2)}]
  \ifnum#7=1 \draw[fine] (0,0) circle[radius=#3]; \else \draw[fine] (90:#3)--(210:#3)--(330:#3)--cycle; \fi
  \foreach \mang/\mc in {90/#4,210/#5,330/#6} {
    \pgfmathsetmacro{\r}{.42*#3*sqrt(abs(\mc))}\pgfmathsetmacro{\pos}{\mc>0?1:0}
    \ifdim\r pt>0.6pt \ifnum\pos=1 \fill[PlateInk] (\mang:#3) circle[radius=\r]; \else \draw[fine,fill=white] (\mang:#3) circle[radius=\r]; \fi
    \else \draw[fine] ($(\mang:#3)+(-.07,0)$)--($(\mang:#3)+(.07,0)$); \fi}
  \end{scope}}
% amplitudes: v1=(1,1,1)/sqrt3; v2=(2,-1,-1)/sqrt6; v3=(0,1,-1)/sqrt2; t(a,b,c)=(c,a,b); b(a,b,c)=(a,c,b)
\def\sA{.5774}\def\sB{.8165}\def\sC{.4082}\def\sD{.7071}

\begin{document}
\begin{tikzpicture}[plate]
\path[use as bounding box] (0,0) rectangle (15,11.6);
% key: the three version sites on the torsion circle
\begin{scope}[shift={(2.5,8.2)}]
  \pgfmathsetmacro{\R}{1.4}
  \draw[fine] (0,0) circle[radius=\R];
  \draw[blift] (0,{-\R-.5})--(0,{\R+.5});
  \draw[tpath] (100:{\R+.28}) arc[start angle=100,end angle=200,radius={\R+.28}];
  \foreach \a in {90,210,330} {\hatom{\R*cos(\a)}{\R*sin(\a)}{.09}{ColDot}}
  \node[lab,anchor=south east] at ($(90:\R)+(-.08,.1)$) {$H$}; \node[lab,anchor=north east] at (210:\R) {$tH$}; \node[lab,anchor=north west] at (330:\R) {$t^2H$};
  \node[lab] at (150:{\R+.62}) {$t$}; \node[lab,anchor=west] at (.1,{\R+.42}) {$b$};
  \node[sub] at (270:{\R+.32}) {$\tau$};
\end{scope}
% the grid
\foreach \col/\x/\a/\b/\c/\ta/\tb/\tc/\ba/\bb/\bc in {1/7.0/\sA/\sA/\sA/\sA/\sA/\sA/\sA/\sA/\sA, 2/10.0/\sB/-\sC/-\sC/-\sC/\sB/-\sC/\sB/-\sC/-\sC, 3/13.0/0/\sD/-\sD/-\sD/0/\sD/0/-\sD/\sD} {
  \modepic{\x}{9.6}{.62}{\a}{\b}{\c}{1}
  \modepic{\x}{7.5}{.62}{\ta}{\tb}{\tc}{1}
  \modepic{\x}{5.4}{.62}{\ba}{\bb}{\bc}{1}
  \node[lab] at (\x,10.55) {$v_\col$};}
\node[lab,anchor=east] at (5.9,9.6) {$v$}; \node[lab,anchor=east] at (5.9,7.5) {$t\,v$}; \node[lab,anchor=east] at (5.9,5.4) {$b\,v$};
\draw[fine] (6.2,11.05)--(6.2,10.9)--(7.8,10.9)--(7.8,11.05); \node[lab,anchor=south] at (7.0,11.07) {$\mathrm A_1$};
\draw[fine] (9.2,11.05)--(9.2,10.9)--(13.8,10.9)--(13.8,11.05); \node[lab,anchor=south] at (11.5,11.07) {$\mathrm E$};
% beneath: the fixed point and the E plane
\begin{scope}[shift={(11.5,2.6)}]
  \pgfmathsetmacro{\L}{1.25}
  \draw[blift] (-\L-.35,0)--(\L+.35,0); \draw[fine] (0,-\L-.25)--(0,\L+.25);
  \node[sub,anchor=south] at (0,\L+.28) {$v_3$};
  \foreach \a/\l/\anc in {0/{v_2}/west,120/{t\,v_2}/south east,240/{t^2v_2}/north east} {\draw[heavy,-{Stealth[length=2mm,width=1.6mm]}] (0,0)--(\a:\L); \node[lab,anchor=\anc] at (\a:{\L+.18}) {$\l$};}
  \draw[tpath] (8:.5) arc[start angle=8,end angle=112,radius=.5];
\end{scope}
\node[sub,anchor=north] at (10.0,.55) {disc area is the amplitude; ink positive, white negative, a dash for zero};
\end{tikzpicture}
\end{document}
